\sectioncenter{Введение}

Многие практические задачи сводятся к поиску наилучшего решения в огромном пространстве вариантов:
подобрать веса модели так, чтобы она верно предсказывала, найти экстремум многоэкстремальной
функции, составить кратчайший маршрут обхода городов. Классические методы оптимизации требуют
гладкости, выпуклости или перебора, который растёт экспоненциально. Вычислительный интеллект
предлагает другой путь: заимствовать механизмы у живой природы. Нервная система учится на примерах,
подстраивая силу связей между нейронами, а биологическая эволюция находит приспособленные организмы
через наследование, изменчивость и отбор. Первый механизм лежит в основе \emph{искусственных нейронных
сетей}, второй~--- в основе \emph{эволюционных вычислений}, самым известным направлением которых являются
\emph{генетические алгоритмы} (ГА).

Обе парадигмы решают по сути задачу оптимизации, но разными способами: нейронная сеть движется по
градиенту функции ошибки, а ГА работает с популяцией решений и использует только значения целевой
функции, поэтому применим к негладким, дискретным и многоэкстремальным задачам. Эти подходы дополняют друг друга: генетическим алгоритмом можно подбирать
веса и архитектуру сети (нейроэволюция), а нейронную сеть использовать как быструю модель
приспособленности внутри ГА~\cite{rutkovskaya2006}.

Пособие состоит из двух частей. \emph{Теоретическая часть} содержит необходимый материал по машинному
обучению и нейронным сетям (разделы~1--4), по генетическим алгоритмам (разделы~5--7) и по современным
безградиентным методам~--- CMA-ES, байесовской оптимизации, NSGA-II и эволюции программ с помощью больших
языковых моделей (раздел~8). Нейросетевая
часть опирается на пособие В.\,А.~Головко и В.\,В.~Краснопрошина~\cite{golovko2017} и классические модели,
на которых исторически построены задания (MATLAB Neural Network Toolbox~\cite{medvedev2002}),
генетическая~--- на курс В.\,Ю.~Скобцова, Н.\,В.~Лапицкой и С.\,Н.~Нестеренкова~\cite{skobtsov2018} и
монографии Д.~Холланда, Д.~Голдберга и З.~Михалевича~\cite{holland1975, goldberg1989, michalewicz1996}.
Подразделы <<Реализация на Python>> показывают, как записать алгоритм на NumPy и с какими функциями
библиотек scikit-learn, PyTorch, MiniSom, SciPy, DEAP, \code{cma}, Optuna и \code{pymoo} сравнить результат.

\emph{Лабораторный практикум} включает девять работ (таблица~\ref{tab:labs}). Каждая содержит цель,
постановку задачи, варианты, требования к программе, порядок выполнения, пример выполнения для
варианта~3 с графиками и таблицами реальных вычислительных экспериментов и контрольные вопросы.

\renewcommand{\thetable}{\arabic{table}}% во введении таблица нумеруется без номера раздела
\begin{table}[!htb]
\caption{Лабораторные работы, разделы теории и инструменты Python}
\label{tab:labs}
\centering\small
\begin{tabularx}{\textwidth}{c>{\raggedright}X>{\centering}p{1.6cm}>{\raggedright\arraybackslash}p{5.2cm}}
\toprule
ЛР & Тема & Теория & Python \\
\midrule
1 & Классификация персептроном & 1, 2 & NumPy; \code{sklearn Perceptron} \\
2 & Прогнозирование линейной сетью & 2 & \code{np.linalg.lstsq}, \code{sliding\_window\_view} \\
3 & Аппроксимация двухслойной сетью & 3 & NumPy; \code{MLPRegressor}, PyTorch, SciPy \\
4 & Кластеризация картой Кохонена & 4 & NumPy; MiniSom, \code{KMeans} \\
5 & Простой ГА: функция одной переменной & 5, 6 & NumPy; \code{minimize\_scalar} \\
6 & Бинарный и вещественный ГА: функция двух переменных & 5, 6 & NumPy; \code{differential\_evolution} \\
7 & Перестановочный ГА: задача коммивояжёра & 6, 7 & NumPy; DEAP, \code{tsplib95} \\
8 & CMA-ES, подбор гиперпараметров, NSGA-II & 8 & NumPy; \code{cma}, Optuna, \code{pymoo} \\
9 & Эволюция эвристик с языковой моделью & 8 & NumPy; Anthropic SDK \\
\bottomrule
\end{tabularx}
\end{table}

Особое внимание уделено \emph{экспериментальной культуре}. ГА и обучение сети со случайной
инициализацией~--- стохастические процедуры, поэтому результат подтверждается серией независимых
запусков (доля успехов, медиана, худший случай) и сравнением с эталоном: аналитическим экстремумом,
пределом точности кодирования или известным оптимальным туром.

\renewcommand{\thetable}{\thesection.\arabic{table}}
